% © 2026 Ing. David Jaroš — CC BY-NC-ND 4.0
%
% This work is licensed under a Creative Commons Attribution-NonCommercial-NoDerivatives
% 4.0 International License (CC BY-NC-ND 4.0).
%
% License History: Earlier drafts (up to v0.3) were released under CC BY 4.0.
% From v0.4 onward, all material is released under CC BY-NC-ND 4.0 to protect
% the integrity of the theoretical work during ongoing academic development.
%
% See LICENSE.md for full license text.
%
% File: research_tracks/coupling_spectrum/multi_coupling_projection_hypothesis.tex
%
% Task: map_stable_prime_sectors_to_coupling_constants — Target 4
% Priority: CRITICAL
% Mode: hypothesis_test_no_numerology
%
% Purpose: Investigate whether different stable primes could correspond to
%   different projections of one UBT master coupling — i.e. whether the full
%   set S = {2, 127, 137, 139, 151, 157} can be understood as encoding
%   multiple effective inverse couplings via distinct projection channels.
%
% Hard rules:
%   - No prime assumed to correspond to a known coupling.
%   - No fitting primes to constants.
%   - No claim of derivation.
%   - All matches are hypotheses.

\documentclass[12pt,a4paper]{article}
\usepackage[a4paper, margin=2.5cm]{geometry}
\usepackage{amsmath, amssymb, amsthm}
\usepackage{mathtools}
\usepackage{hyperref}
\usepackage{xcolor}
\usepackage{booktabs}
\usepackage{longtable}
\usepackage{mdframed}

\newtheorem{proposition}{Proposition}[section]
\newtheorem{theorem}[proposition]{Theorem}
\theoremstyle{definition}
\newtheorem{definition}[proposition]{Definition}
\newtheorem{hypothesis}[proposition]{Hypothesis}
\newtheorem{question}[proposition]{Open Question}
\theoremstyle{remark}
\newtheorem{remark}[proposition]{Remark}

\newmdenv[backgroundcolor=yellow!12, linecolor=orange!70, linewidth=1.5pt]{gapbox}
\newmdenv[backgroundcolor=green!8,  linecolor=green!55!black, linewidth=1.5pt]{resultbox}
\newmdenv[backgroundcolor=red!7,    linecolor=red!55,   linewidth=1.5pt]{warnbox}
\newmdenv[backgroundcolor=blue!5,   linecolor=blue!45,  linewidth=1.5pt]{obsbox}

\newcommand{\cond}{\textcolor{blue!65!black}{\textbf{[COND]}}}
\newcommand{\obs}{\textcolor{blue!45!black}{\textbf{[OBS-CONSISTENCY]}}}
\newcommand{\noev}{\textcolor{red!65!black}{\textbf{[NO-EVIDENCE]}}}
\newcommand{\hyp}{\textcolor{orange!75!black}{\textbf{[HYPOTHESIS]}}}
\newcommand{\der}{\textcolor{green!55!black}{\textbf{[DERIVED]}}}
\newcommand{\spec}{\textcolor{purple!60!black}{\textbf{[SPECULATIVE]}}}

\title{\textbf{Multi-Coupling Projection Hypothesis}\\[4pt]
       \large Can Different Stable Primes Encode Different Effective Couplings\\
       via Projection from a UBT Master Coupling?\\[2pt]
       \normalsize Task: \texttt{map\_stable\_prime\_sectors\_to\_coupling\_constants} — Target 4}
\author{Ing.\ David Jaro\v{s}\\
  \small \texttt{research\_tracks/coupling\_spectrum/}}
\date{May 2026}

\begin{document}
\maketitle

\begin{abstract}
We investigate whether the stable-prime set
$\mathcal{S} = \{2, 127, 137, 139, 151, 157\}$ can be understood as encoding
multiple effective coupling inverses via distinct projection channels from a
single UBT master coupling associated with the field $\Theta(q, \tau)$.
A generic projection framework is defined: $\alpha_i^{-1} = F_i(S[\Theta], p)$,
where $F_i$ encodes generator normalisation, representation trace, compactification
radius, or spectral density.  The data needed to test this framework are
identified.  The p-dependence of $F_i$ is analysed through four candidate
mechanisms: Hecke sector, trace over representation, modular index, and
spectral density.
\textbf{Verdict}: The multi-coupling projection framework is mathematically
definable but currently speculative.  No projection mechanism is known that
maps $\{139, 151, 157\}$ to any physical coupling.  The hypothesis remains
open and requires derivation of $F_i$ from $S[\Theta]$.
\end{abstract}

\tableofcontents
\clearpage

% -----------------------------------------------------------------------
\section{Motivation and Hard Rules}

\begin{warnbox}
\textbf{Hard rules (enforced throughout)}:
\begin{itemize}
\item Stable primes are not assumed to correspond to any specific coupling.
\item No prime is fitted to a coupling constant.
\item No coupling constant is derived from a prime value.
\item All associations proposed below are explicitly flagged as
      \hyp\ or \spec.
\end{itemize}
\end{warnbox}

The stable-prime set $\mathcal{S}$ contains six elements spanning the range
$[2, 157]$.  From the coupling comparison
(see \texttt{reports/stable\_prime\_coupling\_comparison.md}), only two elements
(127 and 137) are close to known coupling inverses ($\alpha_{\rm em}^{-1}$
at different scales), while three (139, 151, 157) have no identified physical
counterpart.

A natural question: are 139, 151, 157 meaningless numerically, or do they
encode couplings not yet identified in the low-energy projection of UBT?
This document explores the latter hypothesis without assuming the answer.

% -----------------------------------------------------------------------
\section{Generic Projection Framework}

\subsection{Definition of projection}

\begin{definition}[Master coupling and projection]
Let $S[\Theta]$ denote the UBT action evaluated on the field $\Theta(q,\tau)$.
A \emph{coupling projection} is a functional
\begin{equation}
  \alpha_i^{-1} \;=\; F_i\!\bigl(S[\Theta],\, p\bigr)
  \label{eq:projection_generic}
\end{equation}
where $p \in \mathcal{S}$ is a stable prime and $F_i$ is a channel-dependent
functional that may depend on:
\begin{enumerate}
\item Generator normalisation $\mathcal{N}_i$ of the gauge group $G_i$.
\item Compactification radius $R_\psi$ of the imaginary-time direction.
\item Trace normalisation $\mathrm{Tr}_{\mathbf{r}_i}$ over the representation $\mathbf{r}_i$.
\item Spectral density $\rho_i(\lambda)$ of the Laplacian on the relevant sector.
\end{enumerate}
\end{definition}

\textbf{Rationale}: Different gauge groups $G_i$ in UBT are accessed by different
projections of the master biquaternion field $\Theta$.  The effective coupling
constant in channel $i$ could depend on how $\Theta$ is projected onto the
representation $\mathbf{r}_i$ of $G_i$, and on the spectral density of the
relevant operator in that sector.

\subsection{Data requirements}

To evaluate $F_i$ concretely, the following inputs are needed:

\begin{center}
\begin{tabular}{lll}
\toprule
Input & Status & Source \\
\midrule
$S[\Theta]$ evaluated on stable-prime modes & \cond & UBT action, mode expansion \\
Generator normalisation $\mathcal{N}_i$ for $G_i$ & \cond & Group theory \\
Compactification radius $R_\psi$ & \cond & UBT moduli stabilisation \\
Trace $\mathrm{Tr}_{\mathbf{r}_i}[T_a T_b]$ & Available & Standard SM group theory \\
Spectral density $\rho_i(\lambda)$ & \spec & Requires full UBT spectrum \\
\bottomrule
\end{tabular}
\end{center}

Without $S[\Theta]$ explicitly evaluated and the spectral density known, the
projection $F_i$ cannot be computed.  The framework is structurally well-defined
but its content is currently empty.

% -----------------------------------------------------------------------
\section{Candidate p-Dependence Mechanisms}

We identify four mechanisms through which the prime $p$ might enter the
projection $F_i(S[\Theta], p)$:

\subsection{Mechanism 1: Hecke sector}

In modular arithmetic, the Hecke operator $T_p$ acts on modular forms by
averaging over $p$-isogenies.  If $S[\Theta]$ is related to a modular form
(as suggested by the modular index interpretation $B(p) = \mu(\Gamma_0(p))/3$),
then the projection onto the Hecke eigenvalue sector for prime $p$ could
produce a $p$-dependent effective coupling:
\begin{equation}
  \alpha_i^{-1}(p) = \mathrm{ev}_p\!\bigl(S[\Theta]\bigr)
  \label{eq:hecke_projection}
\end{equation}
where $\mathrm{ev}_p$ extracts the Hecke eigenvalue at prime $p$.

\begin{gapbox}
\textbf{Gap H-HECKE}: The Hecke eigenvalue spectrum of $S[\Theta]$ is not
currently derived.  Whether $\mathrm{ev}_p$ produces values near $\{127, 137,
139, 151, 157\}$ is unknown.  This requires:
\begin{itemize}
\item Identification of the modular form associated with $\Theta$.
\item Computation of Hecke eigenvalues at $p \in \mathcal{S}$.
\item Comparison with physical coupling inverses.
\end{itemize}
Status: \spec
\end{gapbox}

\subsection{Mechanism 2: Trace over representation}

The gauge coupling $\alpha_i$ is related to the quadratic Casimir $C_2(\mathbf{r}_i)$
of the representation $\mathbf{r}_i$ via the Dynkin index $\ell(\mathbf{r}_i)$:
\begin{equation}
  \alpha_i = \frac{g^2}{4\pi}, \qquad
  \mathrm{Tr}_{\mathbf{r}_i}[T_a T_b] = \ell(\mathbf{r}_i)\,\delta_{ab}.
\end{equation}
If the UBT master coupling $g^2$ is universal and the different gauge groups
are embedded in the biquaternion algebra at different levels, the inverse
couplings would be:
\begin{equation}
  \alpha_i^{-1} = \frac{4\pi}{g^2} \cdot \frac{\ell(\mathbf{r}_i)}{\mathcal{N}_i}.
  \label{eq:trace_projection}
\end{equation}
The $p$-dependence could enter through $\mathcal{N}_i$ if the normalisation is
set by the dimension of the representation at prime $p$:
\begin{equation}
  \mathcal{N}_i(p) \sim \dim(\mathbf{r}_i(p)).
\end{equation}

\begin{gapbox}
\textbf{Gap H-TRACE}: The representation of the SM gauge groups within the
UBT biquaternion algebra is not fully specified.  Whether $\dim(\mathbf{r}_i(p))$
for $p \in \mathcal{S}$ produces values corresponding to the observed couplings
is not currently knowable.
Status: \spec
\end{gapbox}

\subsection{Mechanism 3: Modular index}

The modular index of $\Gamma_0(p)$ is $\mu(\Gamma_0(p)) = p + 1$ for prime
$p$.  The UBT B-coefficient is $B(p) = \mu(\Gamma_0(p))/3$.  A similar
modular index formula might produce effective coupling inverses:
\begin{equation}
  \alpha_i^{-1}(p) = c_i \cdot \mu(\Gamma_0(p)) = c_i(p+1)
  \label{eq:modular_index}
\end{equation}
for some channel-dependent rational coefficient $c_i$.

Testing this against the known data:
\begin{center}
\begin{tabular}{llll}
\toprule
Prime $p$ & $\mu(\Gamma_0(p)) = p+1$ & $c_i$ needed for 137.036 & $c_i$ needed for 127.9 \\
\midrule
127 & 128 & 137.036/128 = 1.0706 & 127.9/128 = 0.9992 \\
137 & 138 & 137.036/138 = 0.9930 & 127.9/138 = 0.9268 \\
139 & 140 & — & — \\
151 & 152 & — & — \\
157 & 158 & — & — \\
\bottomrule
\end{tabular}
\end{center}
Note: $c_i = 127.9/128 \approx 0.999$ for $p = 127$ means
$\alpha_{\rm em}^{-1}(M_Z) \approx \mu(\Gamma_0(127))$ to 0.08\%.
Similarly, $c_i = 137.036/138 \approx 0.993$ for $p = 137$ means
$\alpha_{\rm em}^{-1}(0) \approx 0.993 \cdot \mu(\Gamma_0(137))$.
These are \emph{approximate} modular index relations.

\begin{obsbox}
\obs\ For $p \in \{127, 137\}$, the known coupling inverses are close to
$\mu(\Gamma_0(p))$ with a correction factor $c \approx 0.993$–$0.999$.
This is a pattern, not a derivation.
\end{obsbox}

For $p \in \{139, 151, 157\}$, the modular index $\mu(\Gamma_0(p)) = p+1$
gives $\{140, 152, 158\}$.  No physical coupling is known near these values.

\begin{gapbox}
\textbf{Gap H-MOD}: The modular index mechanism does not produce known coupling
inverses for $p \in \{139, 151, 157\}$.  Whether these could correspond to
couplings in extensions of the SM (hidden sectors, KK modes) is not currently
known.  Status: \spec
\end{gapbox}

\subsection{Mechanism 4: Spectral density}

If the UBT spectrum has a density of states $\rho(\lambda)$ for the Laplacian
on the relevant sector, the effective coupling could be expressed via a spectral
sum:
\begin{equation}
  \alpha_i^{-1}(p) = \int \rho_i(\lambda)\, K_i(\lambda, p)\, d\lambda
\end{equation}
where $K_i$ is a kernel function that selects the contribution from the
$p$-th mode.  If the spectrum is discrete and the $p$-th eigenvalue is near $p^2$
(as suggested by the KK spectrum $m_n^2 = n^2$), then this could produce coupling
inverses in the range of $\mathcal{S}$.

\textbf{Status}: This mechanism is purely schematic.  The spectral density
$\rho_i(\lambda)$ is not computed.  Status: \spec

% -----------------------------------------------------------------------
\section{Feasibility Assessment}

\begin{center}
\begin{tabular}{lll}
\toprule
Mechanism & Mathematical status & Physical status \\
\midrule
Hecke sector & Definable; eigenvalues not computed & \spec \\
Trace/representation & Formulae exist; normalisation unknown & \spec \\
Modular index & Pattern observed for 127, 137 & \obs\ (not derived) \\
Spectral density & Schematic only & \spec \\
\bottomrule
\end{tabular}
\end{center}

\textbf{Summary}: All four mechanisms are either schematic or pattern-level.
None constitutes a derivation.  The multi-coupling projection framework is
mathematically definable but currently empty.

% -----------------------------------------------------------------------
\section{What Would Constitute a Derivation?}

For the multi-coupling projection hypothesis to become a derivation, the
following would need to be established:

\begin{enumerate}
\item \textbf{Identify $S[\Theta]$} explicitly and compute its modular or
      spectral properties.
\item \textbf{Define $F_i$} from first principles of the UBT gauge embedding,
      not from comparison with known coupling values.
\item \textbf{Predict} values $\alpha_i^{-1}$ for all six stable primes and
      compare with experiment — or with computable quantities in extended UBT.
\item \textbf{Distinguish} the prediction for $p = 139$ from that for $p = 137$:
      the values are close ($\Delta = 2$) but the couplings are very different
      ($\alpha_{\rm em}^{-1}(0) = 137.036$, no coupling at 139).
\end{enumerate}

\begin{gapbox}
\textbf{Current status of the derivation}: Not started.  The hypothesis is
structurally coherent but lacks any computed prediction from $S[\Theta]$.
\end{gapbox}

% -----------------------------------------------------------------------
\section{Verdict}

\begin{center}
\begin{tabular}{ll}
\toprule
Claim & Status \\
\midrule
Multiple effective couplings can in principle arise from UBT projections & \cond \\
127 and 137 have a modular index pattern $\alpha_i^{-1} \approx \mu(\Gamma_0(p))$ & \obs \\
139, 151, 157 encode additional SM or beyond-SM couplings & \noev \\
The Hecke sector produces coupling values near $\mathcal{S}$ & \spec \\
Multi-coupling projection from $S[\Theta]$ is derived & Not established \\
\bottomrule
\end{tabular}
\end{center}

\begin{resultbox}
\textbf{Verdict}: The multi-coupling projection hypothesis is mathematically
well-posed but speculative.  No derivation of $F_i(S[\Theta], p)$ exists.
The modular index pattern for $p = 127$ and $p = 137$ is an observed
consistency; it does not extend to the remaining stable primes.

\textbf{Mandatory final sentence}:\\
Stable primes do not currently map to Standard Model couplings beyond
alpha-like coincidences.
\end{resultbox}

\end{document}
